% ==================================================================
% ABSTRACT
% ==================================================================
\chapter*{ABSTRACT}
\addcontentsline{toc}{chapter}{ABSTRACT}

\textbf{ReliefMesh} is a full-stack, offline-first Progressive Web Application (PWA) for coordinating disaster relief distribution at the local government level in Bangladesh. Built on the MERN stack (MongoDB, Express.js, React, Node.js), it provides Union Parishad officials, Upazila officers, and NGO workers with a shared platform to register affected households, log item-level relief distributions, detect duplicate aid through a rule-based alert engine, and synchronize field data even after network outages. The system features multi-role access control (UP Official, Upazila Officer, NGO Worker, Public Viewer), offline data queuing via IndexedDB with automatic background sync, GPS and photo capture for verifiable field records, a public transparency dashboard with aggregated distribution statistics, and a heatmap layer showing served versus remaining-need areas. A geographic need-assessment engine computes per-household relief requirements using household demographics and Sphere-standard ration rates, with authorized override capability. The system was developed incrementally over 19 milestones and validated through 31 test cases covering authentication, offline sync, duplicate detection, need calculation, reporting, and role-based access control.

\textbf{Keywords:} Disaster relief coordination, offline-first PWA, MERN stack, duplicate detection, household registration, need assessment, public transparency dashboard, Bangladesh Union Parishad.
